% !TeX program = xelatex
\documentclass[a4paper,10pt]{article}
\usepackage[top=17mm,bottom=17mm,left=17mm,right=17mm,footskip=9mm]{geometry}
\usepackage{fontspec,xeCJK,amsmath,amssymb,booktabs,tabularx,array,pgfplots,hyperref}
\setmainfont{Times New Roman}
\setCJKmainfont[AutoFakeBold=1.3]{DFKai-SB}
\pgfplotsset{compat=1.18}
\hypersetup{hidelinks,pdftitle={微觀引擎模擬報告},pdfauthor={}}
\setlength{\parindent}{2em}
\setlength{\parskip}{1pt}
\setlength{\tabcolsep}{4pt}
\renewcommand{\arraystretch}{1.2}
\pagestyle{plain}
\renewcommand{\refname}{\normalfont\bfseries\fontsize{12}{16}\selectfont 參考文獻}
\newcommand{\sectiontitle}[2]{\par\addvspace{12pt}{\noindent\bfseries\fontsize{12}{16}\selectfont #1、#2\par}\nobreak\vspace{4pt}}
\newcommand{\captiontext}[1]{\par\vspace{3pt}{\noindent\fontsize{9.5}{13}\selectfont #1\par}}
\makeatletter
\renewenvironment{thebibliography}[1]
 {\par\addvspace{12pt}{\noindent\bfseries\fontsize{11}{15}\selectfont 參考文獻\par}\nobreak\vspace{4pt}
  \list{\@biblabel{\@arabic\c@enumiv}}
  {\settowidth\labelwidth{\@biblabel{#1}}\leftmargin\labelwidth
   \advance\leftmargin\labelsep\usecounter{enumiv}
   \renewcommand\theenumiv{\@arabic\c@enumiv}
   \setlength{\topsep}{1pt}\setlength{\parsep}{0pt}\setlength{\itemsep}{1pt}}\sloppy}
 {\endlist}
\makeatother
% Generated from saved JSON by reports/generate_report.py; do not edit.
\newcommand{\PosMean}{1.2770}
\newcommand{\PosCI}{0.0158}
\newcommand{\PosEta}{0.263}
\newcommand{\RateFrom}{\(0.746\pm0.066\)}
\newcommand{\RateTo}{\(1.366\pm0.015\)}
\newcommand{\RateEtaFrom}{0.196}
\newcommand{\RateEtaTo}{0.275}
\newcommand{\ScaleFrom}{\(1.243\pm0.037\)}
\newcommand{\ScaleTo}{\(1.301\pm0.019\)}
\newcommand{\MoveFrom}{\(-14.732\pm0.343\)}
\newcommand{\MoveTo}{\(0.276\pm0.081\)}
\newcommand{\TempSlope}{0.6857855322222712}
\newcommand{\TempIntercept}{-0.014593010417614405}
\newcommand{\TempSlopeShort}{0.6858}
\newcommand{\TempSE}{0.0040}
\newcommand{\TempB}{-0.0146}
\newcommand{\TempR}{0.99983}
\newcommand{\RatioSlope}{0.9769445526608586}
\newcommand{\RatioIntercept}{0.005238570281702937}
\newcommand{\RatioSlopeShort}{0.9769}
\newcommand{\RatioSE}{0.0036}
\newcommand{\RatioB}{0.0052}
\newcommand{\RatioR}{0.99995}
\newcommand{\WallMean}{1.350}
\newcommand{\WallCI}{0.018}
\newcommand{\WallQS}{1.386}
\newcommand{\PeriodMean}{1.401}
\newcommand{\PeriodCI}{0.018}
\newcommand{\PeriodQS}{1.448}

\begin{document}
\fontsize{11}{15.2}\selectfont
\setlength{\abovedisplayskip}{5pt}
\setlength{\belowdisplayskip}{5pt}
\setlength{\abovedisplayshortskip}{3pt}
\setlength{\belowdisplayshortskip}{4pt}
\begin{center}
{\bfseries\fontsize{16}{21}\selectfont 微觀引擎模擬報告\par}
\vspace{5pt}
熱統計物理一
\end{center}
\sectiontitle{一}{模型與控溫方式}
系統由 \(N\) 個二維、互不作用的粒子組成。粒子受諧和位能 \(V=\frac12\kappa(x-x_0)^2\) 束縛；\(x\) 方向為彈性硬牆，\(y\) 方向為週期邊界。控溫採欠阻尼朗之萬方程\cite{seifert}：
\begin{equation}
\mathrm d\mathbf r=\mathbf v\,\mathrm dt,\qquad
m\,\mathrm d\mathbf v=(-\nabla V-\gamma\mathbf v)\,\mathrm dt
+\sqrt{2\gamma k_{\mathrm B}T_{\mathrm b}}\,\mathrm d\mathbf B.
\label{eq:langevin}
\end{equation}
摩擦與隨機力共同控溫。BAOAB\cite{baoab} 的熱浴子步為 \(\mathbf v^+=c\mathbf v^-+\sqrt{k_{\mathrm B}T_{\mathrm b}(1-c^2)/m}\,\boldsymbol\xi\)，\(c=e^{-\gamma\Delta t/m}\)。各粒子、方向及步的 \(\boldsymbol\xi\) 分量皆為獨立標準高斯數；小步長噪聲尺度為 \(\sqrt{\Delta t}\)。

\textbf{熱浴溫度與粒子估溫不同：}\(T_{\mathrm b}\) 是設定值；由二維動能估得的 \(T_K=K/(Nk_{\mathrm B})\) 會波動，非平衡時還可能包含整體漂移。
\captiontext{共用約化單位與設定：\(m=k_{\mathrm B}=\gamma=1\)、\(\Delta t=0.02\)、\(L_x/L_y=24/12\)。以下以 \(w=W_{\mathrm{out}}/N\) 表示每粒子、每循環的輸出功；除另註外，中心固定於 \(x_0=L_x/2\)。}
\sectiontitle{二}{循環操作與正功結果}
循環依序為：固定 \(\kappa_{\max}\) 換至熱浴 \(\to\) 高溫降低曲率 \(\to\) 固定 \(\kappa_{\min}\) 換至冷浴 \(\to\) 低溫提高曲率。控制量為 \(\lambda=\kappa,x_0\)，廣義力為 \(-\partial V/\partial\lambda\)。控制輸入功為
\begin{equation}
\mathrm dW_{\mathrm{in}}=\sum_{i=1}^{N}\left[\frac{(x_i-x_0)^2}{2}\,\mathrm d\kappa-\kappa(x_i-x_0)\,\mathrm dx_0\right],
\qquad W_{\mathrm{out}}=-W_{\mathrm{in}}.
\label{eq:work}
\end{equation}
功由固定位置的位能差累積，熱由熱浴子步的動能差累積。有限步長下 \(\Delta E=Q_{\mathrm{in}}-W_{\mathrm{out}}+R_{\mathrm{num}}\)，\(R_{\mathrm{num}}\) 另記為數值誤差。

取 \(N=64\)、\(T_{\mathrm H}=3\)、\(T_{\mathrm C}=1\)、\(\kappa_{\min}=1\)、\(\kappa_{\max}=4\)，曲率段與換浴段時間分別為 30、10。20 次獨立重複各暖機 2 圈，再平均 5 圈，得到
\begin{equation}
 w=\PosMean\pm\PosCI\quad\text{（95\% 信賴區間）}.
\label{eq:positive}
\end{equation}
\textbf{信賴區間整段大於零，支持系統能輸出平均正功。}
\sectiontitle{三}{準靜態條件的比較}
準靜態要求操作慢於粒子分布回復平衡的時間。表一比較操作時間、曲率尺度及搬運速率；數值均來自實際模擬，\(\pm\) 表示 \(95\%\) 信賴區間。
\captiontext{表一、參數改變對循環輸出功的影響。}
\noindent{\fontsize{10.5}{14}\selectfont
\begin{tabularx}{\linewidth}{@{}>{\raggedright\arraybackslash}p{.23\linewidth}>{\raggedright\arraybackslash}p{.40\linewidth}>{\raggedright\arraybackslash}X@{}}
\toprule
\textbf{改變的條件}&\textbf{輸出功 \(w\) 的變化}&\textbf{結果解讀}\\
\midrule
曲率段時間\newline \(4\to100\)&\RateFrom{} 至 \RateTo&放慢後接近準靜態值 \(1.386\)；每圈功增加不代表功率增加。\\[3pt]
曲率範圍\newline \(0.5\leftrightarrow2\) 改為 \(2\leftrightarrow8\)&\ScaleFrom{} 至 \ScaleTo&弱阱回復較慢；增大曲率後仍受摩擦限制，不能推論越大越好。\\[3pt]
平均搬運速率\newline \(1.33\to0.10\)&\MoveFrom{} 至 \MoveTo&相對同時序固定阱，功損失由 \(16.005\) 降至 \(0.994\)。\\
\bottomrule
\end{tabularx}\par}
\captiontext{表一皆取 \(N=32\)、20 次獨立重複，各暖機 2 圈後記錄 5 圈，其餘同正功設定。曲率尺度比較固定曲率段 30；搬運固定距離 \(D=4\)、曲率段 30，換浴兼搬運段由 3 增至 40。誤差由每次重複的圈平均計算，不把同一軌跡每圈當成獨立樣本。}
曲率段由 4 放慢至 100 時，效率點估計 \(\eta=\langle W_{\mathrm{out}}\rangle/\langle Q_{\mathrm H}\rangle\) 由 \(\RateEtaFrom\) 增至 \(\RateEtaTo\)；\(Q_{\mathrm H}\) 是兩個高溫段的淨吸熱。本模型無再生器，直接換浴仍不可逆，\textbf{慢速不等同 Carnot 效率}。
\newpage
\sectiontitle{四}{輸出功的線性關係}
對無限域諧和阱，固定中心、曲率操作足夠慢且換浴後充分鬆弛時，
\begin{equation}
 w_{\mathrm{qs}}=\frac{k_{\mathrm B}(T_{\mathrm H}-T_{\mathrm C})}{2}
 \ln\!\left(\frac{\kappa_{\max}}{\kappa_{\min}}\right).
\label{eq:qs}
\end{equation}
因此，\textbf{固定曲率比時，\(w\) 對溫差近似線性；固定溫差時，\(w\) 對曲率比的對數近似線性。}
\pgfplotsset{report/.style={width=\linewidth,height=5.0cm,
 axis line style={black,thin},tick style={black,thin},
 tick label style={font=\fontsize{9}{11}\selectfont},
 label style={font=\fontsize{10}{12}\selectfont},
 title style={font=\fontsize{10}{13}\selectfont},
 ylabel={$w$},ylabel style={rotate=-90},clip=true}}
\par\noindent\begin{minipage}{.48\linewidth}
\centering
\begin{tikzpicture}
\begin{axis}[report,title={(甲)\quad 固定曲率比 \(r=4\)},xmin=-.1,xmax=4.1,ymin=-.12,ymax=3,
 xtick={0,1,2,3,4},ytick={0,1,2,3},xlabel={\(\Delta T=T_{\mathrm H}-T_{\mathrm C}\)}]
\addplot[black!45,solid,thick] table[x=x,y=theory,col sep=comma]{data/temperature.csv};
\addplot[black,dashed,domain=0:4,samples=2]{\TempSlope*x+\TempIntercept};
\addplot[black,only marks,mark=*,mark size=1.7pt,error bars/.cd,y dir=both,y explicit]
 table[x=x,y=mean,y error=ci,col sep=comma]{data/temperature.csv};
\end{axis}
\end{tikzpicture}
\end{minipage}\hfill
\begin{minipage}{.48\linewidth}
\centering
\begin{tikzpicture}
\begin{axis}[report,title={(乙)\quad 固定溫差 \(\Delta T=2\)},xmin=-.05,xmax=1.85,ymin=-.12,ymax=2,
 xtick={0,.5,1,1.5},ytick={0,.5,1,1.5,2},xlabel={\(\ln r=\ln(\kappa_{\max}/\kappa_{\min})\)}]
\addplot[black!45,solid,thick] table[x=x,y=theory,col sep=comma]{data/ratio.csv};
\addplot[black,dashed,domain=0:1.791759469,samples=2]{\RatioSlope*x+\RatioIntercept};
\addplot[black,only marks,mark=*,mark size=1.7pt,error bars/.cd,y dir=both,y explicit]
 table[x=x,y=mean,y error=ci,col sep=comma]{data/ratio.csv};
\end{axis}
\end{tikzpicture}
\end{minipage}
\captiontext{圖一、黑點與誤差棒為模擬平均及 \(95\%\) 信賴區間；黑虛線為含截距的線性擬合，灰實線為模型的準靜態參考。}
\captiontext{條件：硬牆、\(N=32\)、\(T_{\mathrm C}=1\)、\(\kappa_{\min}=1\)，其餘共用參數同前；曲率段 100、換浴段 20。每點 20 次獨立重複，各暖機 2 圈後平均 3 圈。甲圖掃描 \(\Delta T=0\sim4\)；乙圖固定 \(T_{\mathrm H}=3\)，掃描曲率比 \(r=1\sim6\)。}
\captiontext{表二、以 \(w=ax+b\) 擬合的結果；SE 為斜率標準誤。}
\noindent{\fontsize{10.5}{14}\selectfont
\begin{tabularx}{\linewidth}{@{}*{5}{>{\centering\arraybackslash}X}@{}}
\toprule
變因 \(x\)&斜率 \(a\pm\mathrm{SE}\)&理想斜率&截距 \(b\)&\(R^2\)\\
\midrule
\(\Delta T\)&\(\TempSlopeShort\pm\TempSE\)&\(0.6931\)&\(\TempB\)&\(\TempR\)\\
\(\ln r\)&\(\RatioSlopeShort\pm\RatioSE\)&\(1\)&\(\RatioB\)&\(\RatioR\)\\
\bottomrule
\end{tabularx}\par}
兩組資料都支持近似線性。斜率與理想值的差異仍須考慮有限速率與統計誤差，\textbf{不能僅憑 \(R^2\) 高就判定已達準靜態}。系統總功為 \(Nw\)，相應的斜率與截距亦乘上 \(N\)。
\sectiontitle{五}{雙週期邊界的處理}
直接把諧和阱兩端接起來會產生接縫。因此雙週期模式改用平滑週期位能；令 \(\theta=2\pi(x-x_0)/L\)，則
\begin{equation}
 V=\frac{\kappa L^2}{4\pi^2}(1-\cos\theta),\qquad
 F_x=-\frac{\kappa L}{2\pi}\sin\theta.
\label{eq:periodic}
\end{equation}
此位能的阱底曲率仍為 \(\kappa\)。力與控制功皆依此位能計算，跨界時保持一致，座標回捲不額外產生功；準靜態參考由此模型的 \(Z_x=\int_0^L e^{-V/(k_{\mathrm B}T)}\,\mathrm dx\) 計算，不能直接套用式\eqref{eq:qs}。
\captiontext{表三、相同循環設定下的兩種模型比較。}
\noindent{\fontsize{10.5}{14}\selectfont
\begin{tabularx}{\linewidth}{@{}*{3}{>{\centering\arraybackslash}X}@{}}
\toprule
模型&\(w\)（\(95\%\) 信賴區間）&準靜態參考\\
\midrule
硬牆諧和阱&\(\WallMean\pm\WallCI\)&\(\WallQS\)\\
雙週期平滑阱&\(\PeriodMean\pm\PeriodCI\)&\(\PeriodQS\)\\
\bottomrule
\end{tabularx}\par}
\captiontext{表三採 \(L=24\)、\(N=32\)、曲率段 100、換浴段 20，其餘同正功設定；20 次獨立重複，各暖機 2 圈後平均 5 圈。}
\textbf{兩種模型的全域位能不同，不能把功的差異全部歸因於邊界。}模擬參數與原始逐圈記錄均隨程式保存。
\begin{thebibliography}{9}
\fontsize{9.2}{12}\selectfont\raggedright
\setlength{\itemsep}{1pt}
\bibitem{seifert} U. Seifert, \textit{Rep. Prog. Phys.} \textbf{75}, 126001 (2012). \href{https://arxiv.org/abs/1205.4176}{arXiv:1205.4176}.
\bibitem{baoab} B. Leimkuhler and C. Matthews, \textit{J. Chem. Phys.} \textbf{138}, 174102 (2013). \href{https://arxiv.org/abs/1304.3269}{arXiv:1304.3269}.
\end{thebibliography}
\end{document}
